\newpage
\section{Przegląd literatury}
Uczenie maszynowe niejednokrotnie w różnych gałęziach przemysłu jest wykorzystywane do analizy skomplikowanych i nieliniowo powiązanych ze sobą danych. \cite{review_predictive_maintenance_2023} Jednym z powszechnych zastosowań jest analiza przeżywalności komponentów i całych silników lotniczych \cite{saxena2008damage}, dzięki której producenci silników (np. General Electric) mogą oszacować ilość napraw w następującym roku dla każdego silnika, która determinuje przychody z tytułu utrzymywania zdatności do lotu floty klienta (linii lotniczej). Te szacowane przychody następnie są raportowane akcjonariuszom co pośrednio wpływa na wycenę akcji. 

Skutkiem zniszczeń maszyn w przemyśle może być strata finansowa lub uszczerbek na zdrowiu bądź strata życia ludzi. Wczesne wykrywanie uszkodzeń poza minimalizacją kosztów eksploatacji i wymiany urządzeń jest nieodzowne dla bezpiecznego i ciągłego funkcjonowania. Z tego powodu implementacja uczenia maszynowego, a w szczególności autoenkoderów do szacowania RUL (ang. Remaining Usuful Life, tj. dosł. Pozostałego Użytecznego Życia) jest istotna i skutecznie wykorzystywana. \cite{autoencoder_fault_detection_2022}

Autoenkodery wyróżniają się na tle innych narzędzi takich jak Analiza Głównych Składowych (ang. Principal Component Analysis) wyższą jakością rekonstrukcji oraz kompresji (zmniejszania wymiarowości) danych. Kompresja jest kluczowa ze względu na nieustannie rosnący popyt na składowanie danych (ang. Data storage) wraz z podażą, która nie rośnie tak szybko jak popyt. Ponadto zmniejszanie wymiarowości sprzyja wizualizacji oraz komunikacji oraz przyspiesza analizę. \cite{hinton2006reducing}

Informatyczne systemy kwantowe są w stanie modelować zjawiska, które przekraczają modelowanie oparte na fizyce klasycznej. Innymi słowy systemy uczenia maszynowego oparte na rozwiązaniach kwantowych przekraczają możliwościami rozpoznawania wzorców klasyczne systemy. \cite{romero2017quantum}